\section{Introduction}

A succinct non-interactive argument of knowledge (SNARK) lets a
prover convince a verifier that it knows a witness to a computation,
using a short proof that is substantially cheaper to verify than
re-executing the computation. Zero-knowledge variants also hide the
witness. SNARKs thus make it possible to verify computations
performed by an untrusted party~\cite{eurocrypt/Groth16,crypto/Setty20}.
Their main computational cost lies with the prover, whose work can
substantially exceed the computation being certified.

For low-precision computations, part of this overhead comes from
the field used by the proof system. Many efficient SNARKs combine a
polynomial interactive oracle proof (PIOP) with a polynomial
commitment scheme (PCS). The PIOP reduces correctness to polynomial
identities, and the PCS authenticates the polynomial evaluations
used to check them~\cite{crypto/Setty20,eurocrypt/ChenBBZ23}.
The algebraic soundness error depends on the degrees of these
identities relative to the size of the challenge field.
Consequently, the proof may require a large field even when each
witness entry occupies only a few bits.

This difference persists after the witness has been committed.
In sumcheck, for example, random challenges combine entries of the
witness tables into new field elements. Even if the entries lie in
a small subfield, the resulting values generally occupy the full
challenge field~\cite{iacr/BagadDDT25}.
Subsequent rounds multiply and store these full-field values.
The prover therefore pays for the challenge-field representation
throughout much of the PIOP, losing the arithmetic and memory
advantages of the original witness.

Quantized machine-learning inference illustrates the gap.
SmoothQuant supports matrix multiplications with 8-bit weights and
8-bit activations~\cite{icml/XiaoLSWDH23}.
In zero-knowledge machine learning (zkML), Kang et al.~\cite{corr/KangHSS22}
prove DNN inference with 8-bit weights and
activations, and intermediate computations of up to 32 bits,
using a 254-bit prime field.
The precision needed to execute the model is thus much smaller
than the field used to prove its execution. A compact witness
commitment addresses only part of this gap: the relation proof
must also exploit the witness representation~\cite{eurocrypt/DiamondP26,iacr/BagadDDT25}.
This leads to the question we study:
\begin{quote}
\emph{Can a SNARK retain small-field arithmetic for most of the
prover's work while drawing its soundness challenges from a large field?}
\end{quote}

\paragraph{Our approach.}
We present \methodname, a framework of packed arguments for vector
relations over a small finite field \(\Field\).
Let \(N=m\kappa\), and let \(\ExtensionField\) be an extension of
\(\Field\) of degree \(\kappa\).
Our central reduction transforms a length-\(N\) relation over
\(\Field\) into a length-\(m\) relation over \(\ExtensionField\).
To obtain this reduction, we arrange each witness vector as an
\(m\times\kappa\) matrix and encode its rows by polynomials over
\(\Field\). An appropriate polynomial identity captures the
original relation. Evaluating this identity at a random
extension-field point produces the shorter relation.
The prover computes the row-polynomial products over \(\Field\)
and invokes the extension-field proof on only \(m\) entries.
Section~\ref{sec:overview} explains the construction.

We instantiate polynomial packing for inner and Hadamard products,
and extend the resulting arguments to permutation constraints.
The two product reductions share a terminal inner-product proof
over \(\ExtensionField\), but use different polynomial encodings
to preserve the original claim.
For \(|\Field|\geq\kappa\) and \(m\geq\kappa\), both reductions
use \(O(\kappa N)\) base-field operations and \(O(N)\) base-field
coordinates of memory under schoolbook arithmetic and precomputed
bases, excluding commitments and openings.
The corresponding uncompressed extension-field proof uses
\(O(\kappa^2N)\) base-field operations and \(O(\kappa N)\)
coordinates.

The commitment interface must account for a further distinction:
the witness is committed using a fixed field basis, whereas the
relation reduction uses packing weights determined by verifier
challenges. We give an opening reduction that connects these two
representations through base-field coordinate identities.
Combining the packed product, permutation, and opening arguments
gives protocols for a single circuit execution and for multiple
executions of the same circuit.
Theorems~\ref{theorem:single-instance-packing}
and~\ref{theorem:multiple-instance-packing} summarize their
efficiency and soundness.

The framework also supports small fields that are not prime.
For binary arithmetizations, the tower
\(\mathbb F_2\subset\mathbb K=\mathbb F_{2^8}\subset\ExtensionField\)
allows byte-field arithmetic in the working-field PIOP while
retaining a larger challenge field.
Section~\ref{sec:binary-packing} gives the required parameters and
Booleanity constraints, and explains how to use a compatible
code-based binary PCS through the code-switching interface of
Section~\ref{sec:pcs}.

\subsection{Related Work}
\label{sec:related-work}

\paragraph{SNARKs and polynomial IOPs.}
The PIOP--PCS framework separates algebraic verification from
cryptographic authentication.
Spartan~\cite{crypto/Setty20} uses sumcheck to verify rank-one
constraint systems, while HyperPlonk~\cite{eurocrypt/ChenBBZ23}
expresses PLONK-style gate and permutation constraints over the
Boolean hypercube. Both use multilinear polynomial representations.
The prover cost has two corresponding components: computing the
PIOP messages and producing the commitments and evaluation proofs.
This distinction matters for small witnesses, since reducing one
cost need not reduce the other.

\paragraph{ECC-based constructions.}
Elliptic-curve cryptography supports pairing-based SNARKs and
discrete-logarithm-based inner-product arguments.
Groth's pairing-based SNARK~\cite{eurocrypt/Groth16} has proofs
of three group elements and circuit-specific preprocessing.
KZG~\cite{asiacrypt/KateZG10} gives constant-size polynomial
commitments and evaluation proofs.
Sonic~\cite{ccs/MallerBKM19}, Marlin~\cite{eurocrypt/ChiesaHMMVW20},
and PLONK~\cite{iacr/GabizonWC19} provide SNARKs with universal,
updatable structured reference strings.
Bulletproofs~\cite{sp/BunzBBPWM18} gives an inner-product argument
with logarithmic communication and no trusted setup, while
Hyrax~\cite{sp/WahbyTSTW18} provides a multilinear commitment with
communication--verification tradeoffs.
These constructions perform polynomial arithmetic in the scalar
field of the chosen group. Cryptographic security constrains the
field size independently of the witness bit width.

\paragraph{Hash-based constructions.}
Hash-based arguments authenticate error-correcting encodings using
Merkle trees. STARKs~\cite{iacr/Ben-SassonBHR18} use this approach
with FRI~\cite{icalp/Ben-SassonBHR18}, an interactive oracle proof
of proximity for Reed--Solomon codes.
Ligero~\cite{dcc/AmesHIV23} uses interleaved Reed--Solomon codes
to obtain sublinear arguments, and
Brakedown~\cite{crypto/GolovnevLSTW23} obtains linear-time proving
from linear-time encodable codes over sufficiently large fields.
BaseFold~\cite{crypto/ZeilbergerCF24} combines foldable codes with
sumcheck to obtain a field-agnostic multilinear PCS.
WHIR~\cite{eurocrypt/ArnonCFY25} gives constrained Reed--Solomon
proximity tests that support evaluation queries with fast verification.
These schemes allow the arithmetic field to be chosen independently
of a group order, subject to their coding and soundness requirements.

Neither commitment family by itself makes the PIOP bit-width-aware.
ECC-based schemes fix an ambient scalar field; hash-based schemes
permit a wider choice of fields. In either case, challenge-dependent
linear combinations can produce full-field values from short
witness entries. Exploiting short values therefore also requires
structure in the relation proof~\cite{iacr/BagadDDT25}.

\paragraph{Small-field embeddings and Binius.}
Cascudo and Giunta~\cite{iacr/CascudoG21} use reverse
multiplication-friendly embeddings to verify Boolean R1CS through
larger-field linear and row checks.
Binius, due to Diamond and Posen~\cite{eurocrypt/DiamondP25},
adapts Brakedown to multilinear commitments over binary towers
without coefficient-wise embedding overhead, and gives corresponding
product, permutation, and lookup arguments.
Their subsequent ring-switching construction~\cite{eurocrypt/DiamondP26}
reduces multilinear commitments over a ground field to commitments
over an extension field. Its commitment cost matches that of the
underlying PCS on inputs of the same bit size, and its opening
reduction adds linear prover and logarithmic verifier overhead.
The compiler applies to general field extensions; the concrete PCS
uses a binary-field adaptation of BaseFold.
Ring-switching addresses the packed commitment interface.
Using the same representation throughout the circuit PIOP requires
additional relation-specific arguments, as in Binius.

\paragraph{Small-value and small-field sumcheck.}
Bagad, Domb, and Thaler~\cite{iacr/BagadDT24} optimize sumcheck
for input polynomials whose values lie in a small subfield, replacing
extension-field multiplications with additional base-field arithmetic.
Bagad, Dao, Domb, and Thaler~\cite{iacr/BagadDDT25} extend this
approach to short integers and small-subfield values, combine it
with equality-polynomial optimizations, and integrate the prover
into Jolt.
Dao and Thaler~\cite{iacr/DaoT24} exploit binary towers in
zero-check, reducing equality-table computation and generalizing
constraint packing with linearly independent extension-field weights.
Wei et al.~\cite{iacr/WeiWXZDWZ25} give packed sumcheck over small
prime fields using repetition and folding to retain base-field
computation, with an application to TFHE bootstrapping.
These techniques improve the relation proof itself, complementing
improvements to the PCS.

\paragraph{Relation to our work.}
\methodname uses polynomial packing to reduce the number of entries
processed by the extension-field proof.
For both inner and Hadamard products, a polynomial identity links
the length-\(N\) base-field relation to a length-\(N/\kappa\)
extension-field inner product. This leaves the row multiplications
over the base field and reduces the dimension of the subsequent
proof.
Our opening reduction authenticates the challenge-dependent
operands of that proof against a fixed-basis witness commitment.
The relation reduction and its commitment interface are developed
in Sections~\ref{sec:overview} and~\ref{sec:pcs}, respectively.
